\setcounter{chapter}{-1}
\chapter{Introduction}
\pagenumbering{arabic} \setcounter{page}{1}

Assignment 1 for COL334 (Computer Networks) asks every team to map as much of the IIT Delhi
campus network as can honestly be measured from an ordinary student device, without
authentication attempts, SNMP, vulnerability scanners, aggressive scan timing, or scanning of
anyone else's device. This report covers both halves of the assignment: \textbf{Part A:
Cartography}, worth 50 marks, and \textbf{Part B: The Stakeout}, the traffic-capture half,
worth 35 marks. Part A is Chapters~\ref{ch:Background}--\ref{ch:Conclusion}; Part B is
Chapters~\ref{ch:B1}--\ref{ch:B3}.

\section{Team}
This report was prepared jointly by:
\begin{itemize}
    \item Tejasvi Shrivastava (2024CS10582)
    \item Yashvardhan Singh Chaudhary (2024CS10300)
\end{itemize}

\section{Rules of Engagement}
All measurements in this report were taken under the constraints set out in the assignment
handout: no authentication attempts, no SNMP, no exploits or vulnerability scanners, scan
timing no more aggressive than \texttt{-T2}, and no scanning of devices that might belong to a
student or staff member rather than to the infrastructure itself.

\section{Capture Windows}
Part B refers throughout to twelve ten-minute capture windows, six per teammate, labelled
\texttt{D\emph{d}W\emph{w}} for day \emph{d} and time band \emph{w}: W1 is the morning band
($\sim$09:00 IST), W2 the afternoon band ($\sim$15:00 IST), and W3 the evening band
($\sim$21:00 IST). Note that the day labels are collection days, not calendar days: D1W1 and
D2W3 were both recorded on 22-08, which turns out to matter in Chapter~\ref{ch:B1}.

\section{Fixed Targets}
Wherever this report refers to ``the targets'', it means the five fixed destinations specified in
the handout:
\begin{itemize}
    \item \textbf{T1}: our default gateway (found per-device, see Chapter~\ref{ch:Background})
    \item \textbf{T2}: \texttt{www.iitd.ac.in}
    \item \textbf{T3}: \texttt{8.8.8.8}
    \item \textbf{T4}: \texttt{www.iitb.ac.in} ($\sim$1100~km)
    \item \textbf{T5}: \texttt{www.mit.edu} ($\sim$12{,}000~km)
\end{itemize}

\section{Chapter List}
\textbf{Chapter~\ref{ch:Background}, Your Path Out.} Our own device-to-Internet path: interface,
gateway, DNS, and the first hops out of campus, each tagged with evidence.

\textbf{Chapter~\ref{ch:Design}, Campus Topology.} The router hunt (A2.1), the wireless survey
(A2.2), and our contribution to the shared class map (A2.3).

\textbf{Chapter~\ref{ch:Implementation}, The Delay Experiment.} Where RTT time goes (A3.1),
making transmission delay visible via a packet-size sweep (A3.2), and watching a queue fill
(A3.3).

\textbf{Chapter~\ref{ch:Evaluation}, Evidence.} The numbered evidence items (E-1 onward)
that the rest of the report cites.

\textbf{Chapter~\ref{ch:Conclusion}, Findings.} Three specific, checkable findings about the
network, each with an alternative explanation we could not rule out.

\textbf{Chapter~\ref{ch:B1}, Two Days, Two Devices (B1).} Twelve ten-minute captures, six per
teammate across two days and three time bands, compared for RTT to T2 and T3, and for what
each machine sent when we were not asking it to.

\textbf{Chapter~\ref{ch:B2}, The Hunt (B2).} How many anomalies were planted in each of our
captures, a display filter and a frame count for every one of them, and why the looser filters
give the wrong answer.

\textbf{Chapter~\ref{ch:B3}, Argue Against Yourself (B3).} The anomaly that would be hardest to
find at scale, the innocent explanation for it, and what an operator sees that we do not.

\section{A Note on Honesty}
The assignment explicitly rewards the distinction between what we observed directly and what
we are inferring, and rewards an honestly reported ``no difference'' or ``couldn't reach'' over a
tidied-up positive result. We have tried to keep to that standard throughout: flagging
inferred elements, reporting dropped probes as data rather than omitting them, and naming a
second explanation wherever our evidence does not fully rule one out.
